\documentclass[10pt,letterpaper]{article}
\usepackage[margin=0.6in]{geometry}
\usepackage{booktabs,amsmath,array}
\usepackage[T1]{fontenc}
\setlength{\parindent}{0pt}
\setlength{\tabcolsep}{4pt}
\renewcommand{\arraystretch}{1.08}
\begin{document}
\begin{center}
{\large\bfseries Table 1. Predictive relevance of potential experience instruments}\par
\smallskip
\textit{Baseline experience specification}\par
\medskip
\end{center}
\small
\noindent\textbf{Panel A. NLSY79}\par\smallskip
\noindent\begin{tabular*}{\textwidth}{@{\extracolsep{\fill}}lrrrrr@{}}
\toprule
Dependent variable & $E$ & $E^2$ & $E^3$ & $S\times E$ & $A\times E$ \\
\midrule
Joint F statistic & 81.60 & 29.87 & 17.90 & 81.05 & 39.72 \\
Joint test p value & $<$0.001 & $<$0.001 & $<$0.001 & $<$0.001 & $<$0.001 \\
Partial $R^2$ & 0.223 & 0.165 & 0.097 & 0.253 & 0.393 \\
$R^2$ & 0.878 & 0.738 & 0.502 & 0.876 & 0.906 \\
Observations & 30,246 & 30,246 & 30,246 & 30,246 & 30,246 \\
Respondents & 7,631 & 7,631 & 7,631 & 7,631 & 7,631 \\
Excluded instruments & 5 & 5 & 5 & 5 & 5 \\
\bottomrule
\end{tabular*}\par\smallskip
\medskip
\noindent\textbf{Panel B. NLSY97}\par\smallskip
\noindent\begin{tabular*}{\textwidth}{@{\extracolsep{\fill}}lrrrrr@{}}
\toprule
Dependent variable & $E$ & $E^2$ & $E^3$ & $S\times E$ & $A\times E$ \\
\midrule
Joint F statistic & 138.85 & 61.84 & 39.18 & 146.36 & 47.24 \\
Joint test p value & $<$0.001 & $<$0.001 & $<$0.001 & $<$0.001 & $<$0.001 \\
Partial $R^2$ & 0.332 & 0.356 & 0.366 & 0.368 & 0.519 \\
$R^2$ & 0.878 & 0.892 & 0.879 & 0.883 & 0.940 \\
Observations & 11,438 & 11,438 & 11,438 & 11,438 & 11,438 \\
Respondents & 2,889 & 2,889 & 2,889 & 2,889 & 2,889 \\
Excluded instruments & 5 & 5 & 5 & 5 & 5 \\
\bottomrule
\end{tabular*}\par\smallskip
\medskip
{\footnotesize
\noindent\textit{Notes.} E denotes actual experience, P potential experience, S schooling, and A the ASVAB score. Each column is a separate regression containing all five excluded instruments: P, P squared, P cubed, S times P, and A times P.\par
\smallskip\noindent Every regression includes schooling, ASVAB, tenure through its cubic term, schooling times tenure, ASVAB times tenure, and indicators for year, sex, race, and original job number. Schooling, ASVAB, and tenure terms are treated as exogenous. Sampling weights are used and inference is clustered by respondent.\par
\smallskip\noindent The sample retains original job numbers 1 to 3, NLSY79 observations through 2004, and all available NLSY97 years. Eligibility follows the existing pooled wage winsorization at the 1st and 99th percentiles. P is survey year less earliest recorded employment year; it includes nonemployment after entry. It is an observed entry proxy and may include student work. NLSY79 starts coded as week zero have unknown dates before 1978 and are left missing.\par
\smallskip\noindent Requiring usable potential experience reduces the eligible NLSY79 sample from 44,532 to 30,246 observations, and NLSY97 from 11,440 to 11,438. Partial R squared measures the additional explanatory power of the full excluded instrument set, conditional on the included controls. Joint F tests have 5 numerator degrees of freedom and the number of respondents less one denominator degrees of freedom. These tests do not establish exclusion validity or conditional strength for the joint IV system.\par
}
\newpage
\begin{center}
{\large\bfseries Table 2. First stage coefficients on excluded instruments}\par
\smallskip
\textit{Coefficients with respondent clustered standard errors in parentheses}\par
\medskip
\end{center}
\small
\noindent\textbf{Panel A. NLSY79}\par\smallskip
\noindent\begin{tabular*}{\textwidth}{@{\extracolsep{\fill}}lrrrrr@{}}
\toprule
Excluded instrument & $E$ & $E^2$ & $E^3$ & $S\times E$ & $A\times E$ \\
\midrule
$P$ & $-0.0630$ & $-6.4439$ & $18.5543$ & $-7.0044^{**}$ & $0.4480$ \\
 & $(0.1961)$ & $(5.7907)$ & $(177.4009)$ & $(2.8594)$ & $(0.4531)$ \\[2pt]
$P^2$ & $0.0312^{***}$ & $0.7560^{***}$ & $-5.2071$ & $0.3863^{***}$ & $0.0540^{***}$ \\
 & $(0.0071)$ & $(0.2250)$ & $(8.8152)$ & $(0.0937)$ & $(0.0171)$ \\[2pt]
$P^3$ & $-0.0006^{**}$ & $0.0033$ & $0.9714^{**}$ & $-0.0072^{**}$ & $-0.0016^{***}$ \\
 & $(0.0002)$ & $(0.0088)$ & $(0.3810)$ & $(0.0031)$ & $(0.0005)$ \\[2pt]
$S\times P$ & $0.0187$ & $0.1579$ & $-2.9134$ & $0.7589^{***}$ & $-0.0462$ \\
 & $(0.0177)$ & $(0.5189)$ & $(15.1641)$ & $(0.2605)$ & $(0.0409)$ \\[2pt]
$A\times P$ & $0.0505^{***}$ & $1.5157^{**}$ & $45.9944^{*}$ & $0.6368^{***}$ & $0.5890^{***}$ \\
 & $(0.0163)$ & $(0.6751)$ & $(25.3627)$ & $(0.2209)$ & $(0.0604)$ \\[2pt]
\midrule
Joint F statistic & 81.60 & 29.87 & 17.90 & 81.05 & 39.72 \\
Partial $R^2$ & 0.223 & 0.165 & 0.097 & 0.253 & 0.393 \\
Observations & 30,246 & 30,246 & 30,246 & 30,246 & 30,246 \\
Respondents & 7,631 & 7,631 & 7,631 & 7,631 & 7,631 \\
\bottomrule
\end{tabular*}\par\smallskip
\medskip
\noindent\textbf{Panel B. NLSY97}\par\smallskip
\noindent\begin{tabular*}{\textwidth}{@{\extracolsep{\fill}}lrrrrr@{}}
\toprule
Excluded instrument & $E$ & $E^2$ & $E^3$ & $S\times E$ & $A\times E$ \\
\midrule
$P$ & $0.2143$ & $-9.4866^{**}$ & $-195.4017^{*}$ & $-5.6471^{*}$ & $-0.1136$ \\
 & $(0.1715)$ & $(4.6595)$ & $(116.3237)$ & $(2.9718)$ & $(0.3907)$ \\[2pt]
$P^2$ & $0.0300^{***}$ & $1.0459^{***}$ & $5.7715$ & $0.3677^{***}$ & $0.0282$ \\
 & $(0.0076)$ & $(0.2057)$ & $(6.1718)$ & $(0.1015)$ & $(0.0201)$ \\[2pt]
$P^3$ & $-0.0006^{***}$ & $-0.0068$ & $0.5378^{**}$ & $-0.0075^{**}$ & $-0.0009$ \\
 & $(0.0002)$ & $(0.0069)$ & $(0.2202)$ & $(0.0030)$ & $(0.0006)$ \\[2pt]
$S\times P$ & $0.0101$ & $0.3451$ & $9.2768$ & $0.8096^{***}$ & $0.0100$ \\
 & $(0.0118)$ & $(0.3262)$ & $(8.1054)$ & $(0.2086)$ & $(0.0255)$ \\[2pt]
$A\times P$ & $0.0108$ & $0.0086$ & $-3.3578$ & $0.1159$ & $0.6142^{***}$ \\
 & $(0.0121)$ & $(0.3025)$ & $(6.9613)$ & $(0.1742)$ & $(0.0467)$ \\[2pt]
\midrule
Joint F statistic & 138.85 & 61.84 & 39.18 & 146.36 & 47.24 \\
Partial $R^2$ & 0.332 & 0.356 & 0.366 & 0.368 & 0.519 \\
Observations & 11,438 & 11,438 & 11,438 & 11,438 & 11,438 \\
Respondents & 2,889 & 2,889 & 2,889 & 2,889 & 2,889 \\
\bottomrule
\end{tabular*}\par\smallskip
\medskip
{\footnotesize
\noindent\textit{Notes.} E denotes actual experience, P potential experience, S schooling, and A the ASVAB score. Each column is a separate regression containing all five excluded instruments: P, P squared, P cubed, S times P, and A times P.\par
\smallskip\noindent Estimates are reported in original units. Standard errors appear beneath coefficients. *** p $<$ 0.01, ** p $<$ 0.05, * p $<$ 0.10. These stars concern individual coefficients, whereas the joint F statistic tests all five excluded instruments together.\par
\smallskip\noindent Every regression includes schooling, ASVAB, tenure through its cubic term, schooling times tenure, ASVAB times tenure, and indicators for year, sex, race, and original job number. Schooling, ASVAB, and tenure terms are treated as exogenous. Sampling weights are used and inference is clustered by respondent.\par
\smallskip\noindent The samples and entry definition are identical to Table 1. Control coefficients are included in estimation and omitted from the display. No structural wage coefficients are estimated in these tables.\par
}
\newpage
\begin{center}
{\large\bfseries Table 3. Occupation history prediction checks}\par
\smallskip
\textit{Supplementary evidence on predictive relevance}\par
\medskip
\end{center}
\small
\noindent\textbf{Panel A. NLSY79}\par\smallskip
\noindent\begin{tabular*}{\textwidth}{@{\extracolsep{\fill}}lrrrr@{}}
\toprule
Current occupation & History F & History p & Full set F & Full set partial $R^2$ \\
\midrule
Professional & 63.27 & $<$0.001 & 59.56 & 0.422 \\
Managers & 34.13 & $<$0.001 & 33.30 & 0.347 \\
Sales & 11.50 & $<$0.001 & 11.09 & 0.227 \\
Clerical & 76.76 & $<$0.001 & 71.02 & 0.383 \\
Craftsmen & 36.54 & $<$0.001 & 35.37 & 0.326 \\
Operatives and laborers & 59.82 & $<$0.001 & 55.77 & 0.359 \\
\bottomrule
\end{tabular*}\par\smallskip
{\footnotesize Observations: 15,089; respondents: 3,746; history indicators tested: 48; total excluded instruments: 53.\par}
\medskip
\noindent\textbf{Panel B. NLSY97}\par\smallskip
\noindent\begin{tabular*}{\textwidth}{@{\extracolsep{\fill}}lrrrr@{}}
\toprule
Current occupation & History F & History p & Full set F & Full set partial $R^2$ \\
\midrule
Professional & 152.97 & $<$0.001 & 142.63 & 0.711 \\
Managers & 1,384.36 & $<$0.001 & 1,354.73 & 0.653 \\
Sales & 5,081.37 & $<$0.001 & 5,401.60 & 0.709 \\
Clerical & 97.03 & $<$0.001 & 91.05 & 0.692 \\
Craftsmen & 3,420.30 & $<$0.001 & 3,610.77 & 0.802 \\
Operatives and laborers & 318.29 & $<$0.001 & 371.14 & 0.741 \\
\bottomrule
\end{tabular*}\par\smallskip
{\footnotesize Observations: 4,555; respondents: 1,111; history indicators tested: 47; total excluded instruments: 52.\par}
\medskip
{\footnotesize
\noindent\textit{Notes.} Each row is a linear probability regression for the indicated current occupation. Service is the omitted outcome category. Each regression includes the five potential experience terms and the same exogenous controls as Table 1, together with interactions between occupation indicators from the previous two survey waves. Only current main job rows with both lagged occupations observed are included. Histories use the previous two scheduled survey waves, including biennial waves; missing interviews are not skipped.\par
\smallskip\noindent One observed history pair is omitted with the intercept. Empty history pairs are excluded. History F tests only the occupation history block, conditional on the potential experience instruments and included controls. Full set F and partial R squared refer to all excluded instruments together. Standard errors are clustered by respondent and sampling weights are used. Joint test p values below 0.001 are displayed as $<$0.001.\par
\smallskip\noindent These are supplementary prediction checks. They do not implement the full transformed wage model or nonlinear IV estimation in Gibbons et al. (2005). The occupation interactions required by the full thesis wage model are not tested here, and relevance alone does not establish instrument validity.\par
\smallskip\noindent Presentation reference: Gibbons, Katz, Lemieux, and Parent (2005), Journal of Labor Economics, 23(4), Table B3. DOI: 10.1086/491606. Formatting assistance: Codex; numerical estimates are read directly from the saved Stata exports.\par
}
\end{document}
